\section{Association}\label{sec:causalityAssociation}

See \charef{cha:probabilityTheory}: Bayesian Networks are tensor networks and probabilistic queries can be answered by their contractions.
We here provide more theory on query simplification and d-separation, to prepare for the study of intervention mechanisms.


\begin{example}[Markov chain]\label{exa:mcAssociation}

    Markov chains can be parametrized in two directions:

    \stantikzpic{

        \drawtextnode{2}{3}{a) Past-to-future representation:}

% Loop to draw the 4 nodes and their sequential connections
        \foreach \i [
            evaluate=\i as \lab using int(\i-2),
            evaluate=\i as \x using int(\i*5 + 1),
            evaluate=\i as \nextx using int((\i+1)*5 + 1),
            evaluate=\i as \nextidx using int(\i+1)] in {0, 1, 2, 3,4} {

            % Draw the current node
            \ifnum
                \i = 2
                \drawvariablenode{\x}{0}{$\catvariableof{\catenumerator}$}{X\i}
            \fi
            \ifnum
                \i < 2
                \drawvariablenode{\x}{0}{$\catvariableof{\catenumerator\lab}$}{X\i}
            \fi
            \ifnum
                \i > 2
                \drawvariablenode{\x}{0}{$\catvariableof{\catenumerator+\lab}$}{X\i}
            \fi

            % Draw the connecting arrow to the next node / coordinate
            \ifnum
                \i = 0
                \draw[->-] (-4,0) -- (X\i);
            \fi
            \ifnum
                \i=4
                \draw[->-] (X\i) -- (25,0);
            \else
                \draw[->-] (X\i) -- (\nextx,0); % Pre-calculates target x to avoid missing node errors
            \fi
        }

% Draw the final ellipsis node
        \drawtextnode{26}{0}{$\cdots$}
        \drawtextnode{-5}{0}{$\cdots$}

        \begin{scope}[shift={(0,-7)}]

            \drawtextnode{2}{3}{a) Future-to-past representation:}

            \foreach \i [
                evaluate=\i as \lab using int(\i-2),
                evaluate=\i as \x using int(\i*5 + 1),
                evaluate=\i as \nextx using int((\i+1)*5 + 1),
                evaluate=\i as \nextidx using int(\i+1)] in {0, 1, 2, 3,4} {

                % Draw the current node
                \ifnum
                    \i = 2
                    \drawvariablenode{\x}{0}{$\catvariableof{\catenumerator}$}{X\i}
                \fi
                \ifnum
                    \i < 2
                    \drawvariablenode{\x}{0}{$\catvariableof{\catenumerator\lab}$}{X\i}
                \fi
                \ifnum
                    \i > 2
                    \drawvariablenode{\x}{0}{$\catvariableof{\catenumerator+\lab}$}{X\i}
                \fi

                % Draw the connecting arrow to the next node / coordinate
                \ifnum
                    \i = 0
                    \draw[->-] (X\i) -- (-4,0);
                \fi
                \ifnum
                    \i=4
                    \draw[->-] (25,0) -- (X\i);
                \else
                    \draw[->-] (\nextx,0) -- (X\i); % Pre-calculates target x to avoid missing node errors
                \fi
            }

            \drawtextnode{26}{0}{$\cdots$}
            \drawtextnode{-5}{0}{$\cdots$}
        \end{scope}
    }

\end{example}

\subsection{Simplification using directed contractions}

\begin{lemma}\label{lem:bnSubContractionAsCondProb}
    Let $\shortcatvariables$ be an enumeration of the variables respecting a topological order on the directed acyclic hypergraph $\graph$.
    For any $\seccatenumerator<\catenumerator$ we build the sets $[\catenumerator]\backslash[\seccatenumerator] = \{\seccatenumerator,\seccatenumerator+1,\ldots,\catenumerator-1\}$ and $\arbset=\left(\bigcup_{l\in[\catenumerator]\backslash[\seccatenumerator]}\parentsof{l}\right)\cap[\seccatenumerator]$.
    Then
    \begin{align*}
        \condprobat{\catvariableof{[\catenumerator]\backslash[\seccatenumerator]}}{\catvariableof{\arbset}}
        = \contractionof{
            \{\condprobat{\catvariableof{l}}{\catvariableof{\parentsof{l}}} \wcols l\in[\catenumerator]\backslash[\seccatenumerator]\}
        }{
            \catvariableof{[\catenumerator]\backslash[\seccatenumerator]},\catvariableof{\arbset}
        }
    \end{align*}
\end{lemma}
\begin{proof}
    By construction we have for $l\in[\catenumerator]\backslash[\seccatenumerator]$ the independencies
    \begin{align*}
        \condindependent{\catvariableof{l}}{\catvariableof{[l]\backslash([\seccatenumerator]\cup\parentsof{l})},\catvariableof{\arbset\backslash\parentsof{l}}}{\catvariableof{\parentsof{l}}}
    \end{align*}
    and therefore
    \begin{align*}
        \condprobat{\catvariableof{l}}{\catvariableof{\parentsof{l}}}
        \otimes \onesat{\catvariableof{[l]\backslash([\seccatenumerator]\cup\parentsof{l})}}
        \otimes \onesat{\catvariableof{\arbset\backslash\parentsof{l}}}
        =  \condprobat{\catvariableof{l}}{\catvariableof{([l]\backslash[\seccatenumerator])\cup\arbsetof{l}}} \, .
    \end{align*}
    It follows that
    \begin{align*}
        &\contractionof{
            \{\condprobat{\catvariableof{l}}{\catvariableof{\parentsof{l}}} \wcols l\in[\catenumerator]\backslash[\seccatenumerator]\}
        }{
            \catvariableof{[\catenumerator]\backslash[\seccatenumerator]},\catvariableof{\arbset}
        }
        %&\quad
        =\contractionof{
            \{\condprobat{\catvariableof{l}}{\catvariableof{([l]\backslash[\seccatenumerator])\cup\arbsetof{l}}} \wcols l\in[\catenumerator]\backslash[\seccatenumerator]\}
        }{
            \catvariableof{[\catenumerator]\backslash[\seccatenumerator]},\catvariableof{\arbset}
        }
        %&\quad
        =\condprobat{\catvariableof{[\catenumerator]\backslash[\seccatenumerator]}}{\catvariableof{\arbset}}
    \end{align*}
    Here we used in the last equation the chain decomposition on the conditioned distribution $\condprobat{\catvariableof{[\catenumerator]\backslash[\seccatenumerator]}}{\catvariableof{\arbset}}$.
\end{proof}

\subsection{Statistical influence}

Let us investigate, how the directionality in a Bayesian network hypergraph can be interpreted as statistical influence between variables.
To be more precise, we study which cores need to be contracted to answer probabilistic queries, and in such way contribute a statistical influence to a query.

% Simplification rules
When conditioning on some variables, and marginalizing others, we can simplify the Bayesian network.
There are two simplification mechanisms:
\begin{itemize}
    \item[(i)] When having a basis vector to a closed variable, we can remove the variable and replace all tensors affected by the variable with contractions with a copy of the basis vector.
    That is, we apply the contraction identity:
    \stantikzpic{
        \drawsquaretensor{0}{1}{$\onehotmapof{\catindex}$}
        \drawddwire{0}{-2}{$\catvariable$}
        \drawvariabledot{0}{-2}
        \draw[->-] (0,-2) to [bend right=20] (-3,-4);
        \drawtextnode{0}{-3.5}{$\cdots$}
        \draw[->-] (0,-2) to [bend right=-20] (3,-4);

        \drawtextnode{5}{-1}{${=}$}

        \drawsquaretensor{10}{-1}{$\onehotmapof{\catindex}$}
        \drawddwire{10}{-4}{$\catvariable$}

        \drawtextnode{13}{-2.5}{$\cdots$}
        \drawsquaretensor{16}{-1}{$\onehotmapof{\catindex}$}
        \drawddwire{16}{-4}{$\catvariable$}
    }
    \item[(ii)] When a closed variables appears only as an outgoing variable, we can replace the corresponding tensor with the trivial tensor.
    This corresponds with the contraction identity:
    \stantikzpic{
        \drawvariabledot{0}{3}
        \drawtdwire{0}{3}{$\catvariableof{\outsymbol}$}
        \drawtensorblock{0}{0}{2.5}{$\condprobat{\catvariableof{\outsymbol}}{\catvariableof{\insymbol}}$}
        \drawtdwire{-1.5}{-1}{}
        \drawtextnode{0}{-2}{\colorlabelsize $\catvariableof{\insymbol}$}
        \drawtdwire{1.5}{-1}{}

        \drawtextnode{5}{0}{${=}$}

        \drawtensorblock{10}{0}{2}{$\ones$}
        \drawtdwire{8.5}{-1}{}
        \drawtextnode{10}{-2}{\colorlabelsize $\catvariableof{\insymbol}$}
        \drawtdwire{11.5}{-1}{}
    }
\end{itemize}

When having a path through the directed hypergraph (which itself does not necessarily respect the directionality), we investigate whether it will be a path in the simplified hypergraph.
Consiering a node $\catvariableof{1}$ and its neighbors $\catvariableof{0}$, $\catvariableof{1}$ in a path, there are three different situations \cite{pearl_causality_2009}:
\begin{itemize}
    \item {\bf Chains:} $\catvariableof{1}$ appears as outgoing and incoming variable.
    \stantikzpic{
        \drawvariablenode{0}{0}{$\catvariableof{0}$}{X0}
        \drawvariablenode{5}{0}{$\catvariableof{1}$}{X1}
        \drawvariablenode{10}{0}{$\catvariableof{2}$}{X2}
        \draw[->-] (X0) -- (X1);
        \draw[->-] (X1) -- (X2);
    }
    When conditioned on $\catvariableof{1}$, the path closes by rule $(i)$.
    \item {\bf Forks:} $\catvariableof{1}$ appears only as outgoing variable.
    \stantikzpic{
        \drawvariablenode{0}{0}{$\catvariableof{0}$}{X0}
        \drawvariablenode{5}{0}{$\catvariableof{1}$}{X1}
        \drawvariablenode{10}{0}{$\catvariableof{2}$}{X2}
        \draw[-<-] (X0) -- (X1);
        \draw[->-] (X1) -- (X2);
    }
    When conditioned on $\catvariableof{1}$, the path also closes by rule $(i)$.
    \item {\bf Colliders:} $\catvariableof{1}$ appears only as incoming variable. % Also (v-structure)
    \stantikzpic{
        \drawvariablenode{0}{0}{$\catvariableof{0}$}{X0}
        \drawvariablenode{5}{0}{$\catvariableof{1}$}{X1}
        \drawvariablenode{10}{0}{$\catvariableof{2}$}{X2}
        \draw[->-] (X0) -- (X1);
        \draw[-<-] (X1) -- (X2);
    }
    When $\catvariableof{1}$ is marginalized out and not conditioned on, simplification rule $(ii)$ disconnects the path.
    When $\catvariableof{1}$ is conditioned on, the path is still active by the tensor $\condprobat{\indexedcatvariableof{1}}{\catvariableof{0},\catvariableof{2},\catvariableof{\parentsof{1}/\{0,2\}}}$
\end{itemize}

We use this intuition to define active trails and d-separation, oriented on \cite[Def.~3.6 and Def.~3.7]{koller_probabilistic_2009}.

\begin{definition}[Active trails and d-separation]
    A trail in a directed hypergraph $\graph=(\nodes,\edges)$ is a enumerated subset of variables $\catvariableof{[r]}$ such that for each $s\in[r-1]$ there is a hyperedge $\edge=(\incomingnodes,\outgoingnodes)$ with $\catvariableof{s}\in\incomingnodes$ and $\catvariableof{s}\in\outgoingnodes$ or $\catvariableof{s}\in\outgoingnodes$ and $\catvariableof{s}\in\incomingnodes$.
    We say that a trail in $\graph$ is blocked giving a set $\nodesc$, if any collider has a node in $\nodesc$ or in $\mathrm{Desc}(\nodesc)$ in the middle, and no chain or fork has a node in $\nodesc$ in the middle.
    Two subsets $\nodesa$ and $\nodesb$ are said to be d-separated given $\nodesc$, if any trail from $\nodesa$ to $\nodesb$ is blocked by $\nodesc$.
    %   the trail contains a chain or fork with a variable in $\catvariableof{\arbsetof{\mathrm{cond}}}$ in the middle, or contains a collider with a variable in $\catvariableof{\arbsetof{\mathrm{cond}}}$ in the middle.
\end{definition}

To show the utility of this definition, let us consider two sets of variables $\nodesa,\nodesb$ d-separated by $~3.\nodesc$.

\begin{theorem}\label{the:dSeparation}
    If and only if $\nodesa,\nodesb$ d-separated by $\nodesc$ then for any Bayesian network on $\graph$
    \begin{align*}
        \condindependent{\catvariableof{\nodesa}}{\catvariableof{\nodesb}}{\catvariableof{\nodesc}} \, .
    \end{align*}
\end{theorem}
\begin{proof}
    Let us first show, that whenever we have $\nodesa,\nodesb$ which are d-separated by $\nodesc$ and a Bayesian network on $\graph$, then $\condindependent{\catvariableof{\nodesa}}{\catvariableof{\nodesb}}{\catvariableof{\nodesc}}$ with respect to that Bayesian network.
    By performing the simplification rules iteratively on the contraction
    \begin{align*}
        \condprobat{\catvariableof{\nodesa\cup\nodesb}}{\indexedcatvariableof{\nodesc}}
    \end{align*}
    and observing that the tensor network decomposes into a tensor product
    \begin{align*}
        \condprobat{\catvariableof{\nodesa\cup\nodesb}}{\indexedcatvariableof{\nodesc}}
        = \condprobat{\catvariableof{\nodesa}}{\indexedcatvariableof{\nodesc}}
        \otimes \condprobat{\catvariableof{\nodesb}}{\indexedcatvariableof{\nodesc}} \, .
    \end{align*}

    To show the contrary, it is enough construct a Bayesian network not satisfying this independence statement once $\nodesa,\nodesb$ are not d-separated by $\nodesc$.
    We refer to Thm.~3.4 in \cite{koller_probabilistic_2009} for such a construction.
\end{proof}

We already characterized in \charef{cha:probabilityTheory} Bayesian networks on a hypergraph as those distributions satisfying the independencies
\begin{align*}
    \condindependent{\catvariableof{\catenumerator}}{\catvariableof{[\catenumerator]\backslash\parentsof{\catenumerator}}}{\catvariableof{\parentsof{\catenumerator}}}
\end{align*}
with respect to a topological order on $\graph$.


%such that a set $\catvariableof{\arbsetof{\mathrm{cond}}}$ blocks every trail from $\catvariableof{0}$ to $\catvariableof{1}$.
%We then consider the distribution
%\begin{align*}
%    \condprobat{\catvariableof{0},\catvariableof{1}}{\catvariableof{\arbsetof{\mathrm{cond}}}} \, .
%\end{align*}
%In the tensor network perspective, after simplification, we observe that this is a contraction of $\condprobat{\catvariableof{0}}{\catvariableof{\arbsetof{\mathrm{cond}}}}$ and $\condprobat{\catvariableof{0}}{\catvariableof{\arbsetof{\mathrm{cond}}}}$, and therefore
%\begin{align*}
%    \condindependent{\catvariableof{0}}{\catvariableof{1}}{\catvariableof{\arbsetof{\mathrm{cond}}}} \, .
%\end{align*}